% !TEX program = lualatex
% =====================================================================
% Defensa de tesis, DECK v4 FORMAL (20 minutos), estilo Metropolis oscuro.
% Universidad Tecnologica de Pereira, MISC, 2026.
% Autores: Alejandro Gomez Huertas y Juan Diego Garzon Ovalle.
% Compilar con LuaLaTeX (dos veces): lualatex beamer_defensa_v4.tex
%
% Diferencias frente a v3:
%   - Intro y marco conceptual condensados a 4 laminas con figuras de resumen.
%   - Las dos graficas de barras (Resultado 1 y Resultado 4) se sustituyen por
%     figuras conceptuales dibujadas en TikZ que representan la tecnica y el concepto.
%   - Recorte a una defensa de 20 minutos. El respaldo conserva el detalle numerico.
% =====================================================================
\documentclass[aspectratio=169,11pt]{beamer}
\usepackage{fontspec}
\usepackage[spanish,es-noshorthands]{babel}
\usepackage{xcolor}
\usepackage{graphicx}
\usepackage{booktabs}
\usepackage{tikz}
\usetikzlibrary{arrows.meta,positioning,fit,backgrounds}
\graphicspath{{../tesis_latex/}{fig/}}

\usetheme{metropolis}
\metroset{sectionpage=progressbar, progressbar=frametitle, numbering=fraction, block=fill}

% -------------------------------------------------- paleta (modo oscuro)
\definecolor{darkbg}{rgb}{0.05,0.10,0.125}
\definecolor{card}{rgb}{0.10,0.15,0.185}
\definecolor{greytx}{rgb}{0.87,0.88,0.90}
\definecolor{amber}{rgb}{1.0,0.80,0.25}
\definecolor{orange}{rgb}{1.0,0.55,0.40}
\definecolor{teal}{rgb}{0.40,0.80,0.92}
\definecolor{seafoam}{rgb}{0.45,0.86,0.70}
\definecolor{coral}{rgb}{1.0,0.46,0.42}
\definecolor{muted}{rgb}{0.62,0.68,0.78}
\definecolor{navy}{rgb}{0.62,0.80,1.0}

% -------------------------------------------------- colores del tema
\setbeamercolor{normal text}{fg=greytx,bg=darkbg}
\setbeamercolor{background canvas}{bg=darkbg}
\setbeamercolor{alerted text}{fg=amber}
\setbeamercolor{example text}{fg=seafoam}
\setbeamercolor{frametitle}{fg=amber,bg=darkbg}
\setbeamercolor{progress bar}{fg=amber,bg=card}
\setbeamercolor{title separator}{fg=amber}
\setbeamercolor{progress bar in section page}{fg=amber,bg=card}
\setbeamercolor{itemize item}{fg=amber}
\setbeamercolor{itemize subitem}{fg=teal}
\setbeamercolor{enumerate item}{fg=amber}
\setbeamercolor{standout}{fg=amber,bg=darkbg}
\setbeamercolor{block title}{fg=darkbg,bg=amber}
\setbeamercolor{block body}{fg=greytx,bg=card}
\setbeamercolor{block title alerted}{fg=white,bg=coral}
\setbeamercolor{block body alerted}{fg=greytx,bg=card}
\setbeamercolor{block title example}{fg=darkbg,bg=seafoam}
\setbeamercolor{block body example}{fg=greytx,bg=card}

% -------------------------------------------------- helpers
\newcommand{\figcard}[2][\linewidth]{\tikz\node[fill=white,rounded corners=2pt,inner sep=2.5pt]{\includegraphics[width=#1]{#2}};}
\newcounter{figc}
\newcounter{tabc}
\newcommand{\figcap}[1]{\stepcounter{figc}\par\vspace{2pt}{\scriptsize\color{muted}\textbf{Figura~\thefigc.}~#1}}
\newcommand{\tabcap}[1]{\stepcounter{tabc}{\scriptsize\color{muted}\textbf{Tabla~\thetabc.}~#1\par}\vspace{3pt}}

% estilos comunes de nodos para las figuras conceptuales
\tikzset{
  gnode/.style={circle,draw=muted,fill=card,minimum size=6.5mm,inner sep=0,font=\scriptsize,text=greytx},
  gill/.style={circle,draw=coral,fill=coral!30,minimum size=6.5mm,inner sep=0,font=\scriptsize,text=white,line width=0.9pt},
  gbox/.style={rounded corners=3pt,draw=teal,fill=card,text=greytx,align=center,inner sep=5pt,font=\footnotesize},
  gflow/.style={-{Stealth[length=2.4mm]},line width=1pt,teal}
}

% punto con barra de error horizontal para la figura de particion (Resultado 1)
% args: x, err, y, color, nombre, valor
\newcommand{\ptmark}[6]{%
  \draw[#4,line width=1.2pt] (#1-#2,#3) -- (#1+#2,#3);
  \draw[#4,line width=1.2pt] (#1-#2,#3-0.12) -- (#1-#2,#3+0.12);
  \draw[#4,line width=1.2pt] (#1+#2,#3-0.12) -- (#1+#2,#3+0.12);
  \node[circle,fill=#4,draw=darkbg,line width=0.4pt,inner sep=2.2pt] at (#1,#3) {};
  \node[#4,font=\scriptsize,anchor=south] at (#1,#3+0.16) {\textbf{#5}~#6};%
}

% -------------------------------------------------- datos del titulo
\title{Estabilidad de Métodos de Explicabilidad (XAI) en GNNs}
\subtitle{Detección de Lavado de Dinero bajo Desbalance Extremo}
\author{Alejandro Gómez Huertas \and Juan Diego Garzón Ovalle}
\institute{{\color{amber}\textsc{Universidad Tecnológica de Pereira}}\\ Facultad de Ingenierías \textperiodcentered\ Maestría en Ingeniería en Sistemas y Computación\\ Director: Ph.D. Cristian Rosero Arias}
\date{Trabajo de grado para optar al título de Magíster \textperiodcentered\ Pereira, Colombia \textperiodcentered\ 2026}

\begin{document}

\maketitle

\begin{frame}{Contenido de la sustentación}
\begin{columns}[T]
\begin{column}{0.62\textwidth}
\begin{enumerate}
  \item[\textcolor{amber}{\textbf{1}}] \textbf{Problema y marco conceptual}
  \item[\textcolor{amber}{\textbf{2}}] \textbf{Metodología: dos ejes}
  \item[\textcolor{amber}{\textbf{3}}] \textbf{Resultados}
  \item[\textcolor{amber}{\textbf{4}}] \textbf{Conclusiones}
\end{enumerate}
\end{column}
\begin{column}{0.38\textwidth}
\begin{block}{Formato}
\small Defensa formal de \alert{20 minutos}. El detalle numérico queda en el respaldo, para las preguntas del jurado.
\end{block}
\end{column}
\end{columns}
\end{frame}

% ##################################################################### 1
\section{Problema y marco conceptual}

% -------- Lamina 1 de 4: motivacion en una figura
\begin{frame}{El lavado es un patrón de red, y la GNN es una caja negra}
\begin{columns}[T]
\begin{column}{0.46\textwidth}
\begin{itemize}\footnotesize
  \item El lavado mueve \alert{2 a 5\% del PIB mundial} y se delata por el \alert{patrón de conexiones}, no por transacciones aisladas.
  \item Las reglas fijas dan \alert{95 a 98\% de falsos positivos}. Las \textbf{GNN} modelan la red y superan a lo tabular, pero son \textbf{cajas negras}.
  \item En un entorno regulado el analista debe \textbf{justificar} cada alerta, y de forma \textbf{reproducible}.
\end{itemize}
\vspace{0.1em}
\begin{block}{}
\footnotesize Si la explicación \textbf{cambia} al recalcular, no sirve para auditar. Por eso la \alert{estabilidad} es la propiedad crítica de esta tesis.
\end{block}
\end{column}
\begin{column}{0.54\textwidth}\centering
\begin{tikzpicture}[scale=0.68, every node/.style={transform shape}]
  % patron de fan-in ilicito
  \node[gill] (h) at (0,0) {};
  \node[gnode] (a) at (-2.1,1.15) {};
  \node[gnode] (b) at (-2.35,0.15) {};
  \node[gnode] (c) at (-2.1,-0.9) {};
  \node[gnode] (n1) at (-0.6,1.5) {};
  \node[gnode] (n2) at (-3.1,-0.7) {};
  \draw[coral,line width=1.6pt] (a)--(h);
  \draw[coral,line width=1.6pt] (b)--(h);
  \draw[coral,line width=1.6pt] (c)--(h);
  \draw[muted,line width=0.6pt] (a)--(n1);
  \draw[muted,line width=0.6pt] (b)--(n2);
  \node[coral,font=\scriptsize] at (-1.15,-1.6) {patrón de lavado};
  % caja GNN
  \node[gbox,draw=amber,minimum width=1.9cm,minimum height=1.0cm] (gnn) at (2.6,0.3) {GNN\\{\scriptsize caja negra}};
  \draw[gflow,teal] (h) -- (gnn);
  % salida
  \node[align=left,font=\footnotesize] (out) at (2.6,-1.55) {\textcolor{coral}{\textbf{Alerta: ilícita}}\\[1pt]\textcolor{amber}{¿por qué?}\\[1pt]\textcolor{amber}{¿sale igual al repetir?}};
  \draw[gflow,amber] (gnn) -- (out);
\end{tikzpicture}
\figcap{De un patrón de conexiones a una alerta sin explicación auditable. \emph{Fuente: elaboración propia.}}
\end{column}
\end{columns}
\end{frame}

% -------- Lamina 2 de 4: pregunta, objetivos, hipotesis
\begin{frame}{Pregunta, objetivos e hipótesis}
\begin{block}{Pregunta de investigación}
\footnotesize ¿Cómo se comporta la \alert{estabilidad} de los métodos XAI sobre GNNs para detección de lavado bajo desbalance, y qué combinación de \textbf{arquitectura}, \textbf{explicador} y \textbf{balanceo} produce la interpretación más robusta y auditable?
\end{block}
\vspace{0.1em}
\begin{columns}[T]
\begin{column}{0.5\textwidth}
{\scriptsize\color{navy}\textbf{Objetivos}}
\begin{itemize}\scriptsize
  \item[\textcolor{amber}{O1}] ¿Se degrada la estabilidad al agravarse el desbalance?
  \item[\textcolor{amber}{O2}] ¿Qué tan resilientes son las 4 arquitecturas?
  \item[\textcolor{amber}{O3}] ¿Influyen las estrategias de balanceo?
  \item[\textcolor{amber}{O4}] Condensar todo en una matriz de recomendación.
\end{itemize}
\end{column}
\begin{column}{0.5\textwidth}
{\scriptsize\color{navy}\textbf{Hipótesis, escritas antes de ver los datos}}
\begin{itemize}\scriptsize
  \item[\textcolor{coral}{H1}] La estabilidad se degradará con el desbalance.
  \item[\textcolor{coral}{H2}] TAGCN será la más estable, por su alcance multi-salto.
  \item[\textcolor{coral}{H3}] Un explicador más estable señalará mejor el patrón real.
\end{itemize}
\end{column}
\end{columns}
\vspace{0.3em}
{\scriptsize\color{muted} Comprometerse con predicciones concretas antes de ver los datos es lo que hace que refutarlas signifique algo.}
\end{frame}

% -------- Lamina 3 de 4: que es explicar una GNN (pipeline conceptual)
\begin{frame}{Qué es explicar una GNN}
\begin{columns}[T]
\begin{column}{0.4\textwidth}
\footnotesize
Una GNN actualiza cada nodo \textbf{agregando a sus vecinos}, capa por capa. La predicción depende del \textbf{vecindario}, no solo del nodo.
\vspace{0.4em}

\textbf{\color{navy}Cuatro arquitecturas} por su forma de agregar: GCN promedia, GraphSAGE muestrea, GAT usa atención, TAGCN alcanza varios saltos.
\vspace{0.4em}

\textbf{\color{navy}Tres explicadores} post-hoc sobre el modelo ya entrenado: GNNExplainer, PGExplainer y GNNShap.
\vspace{0.3em}

\begin{block}{}
\footnotesize Explicar es señalar \textbf{qué aristas} y \textbf{qué features} sostuvieron la predicción.
\end{block}
\end{column}
\begin{column}{0.6\textwidth}\centering
\begin{tikzpicture}[scale=0.92, every node/.style={transform shape}]
  % paso de mensajes
  \node[gnode,draw=amber,line width=1pt] (v) at (0,0) {$v$};
  \node[gnode] (m1) at (-1.35,0.85) {};
  \node[gnode] (m2) at (-1.55,-0.15) {};
  \node[gnode] (m3) at (-1.15,-1.0) {};
  \draw[gflow] (m1)--(v);
  \draw[gflow] (m2)--(v);
  \draw[gflow] (m3)--(v);
  \node[muted,font=\scriptsize] at (-1.35,-1.6) {paso de mensajes};
  % modelo
  \node[gbox,draw=amber] (mod) at (2.15,0) {GNN entrenada\\{\scriptsize GCN, GraphSAGE,}\\[-1pt]{\scriptsize GAT, TAGCN}};
  \draw[gflow,amber] (v) -- (mod);
  % explicador
  \node[gbox] (exp) at (5.0,0) {Explicador\\{\scriptsize GNNExplainer,}\\[-1pt]{\scriptsize PGExplainer, GNNShap}};
  \draw[gflow,amber] (mod) -- (exp);
  % salidas
  \node[gbox,draw=seafoam,minimum width=2.2cm] (mask) at (5.0,-1.95) {máscara de aristas};
  \node[gbox,draw=seafoam,minimum width=2.2cm] (rank) at (2.1,-1.95) {ranking de features\\{\scriptsize 1, 2, 3, \ldots\ de 166}};
  \draw[gflow,seafoam] (exp) -- (mask);
  \draw[gflow,seafoam] (exp.south) to[out=-120,in=20] (rank.east);
\end{tikzpicture}
\figcap{Del paso de mensajes a las dos salidas del explicador. \emph{Fuente: elaboración propia.}}
\end{column}
\end{columns}
\end{frame}

% -------- Lamina 4 de 4: tres propiedades independientes
\begin{frame}{Tres propiedades que no son lo mismo}
\begin{columns}[T]
\begin{column}{0.32\textwidth}
\begin{block}{Estabilidad}\small ¿La explicación se reproduce entre semillas y perturbaciones?\end{block}
\end{column}
\begin{column}{0.32\textwidth}
\begin{block}{Plausibilidad}\small ¿Señala el patrón real de lavado (\emph{ground-truth})?\end{block}
\end{column}
\begin{column}{0.32\textwidth}
\begin{block}{Fidelidad}\small ¿Refleja de verdad la decisión interna del modelo?\end{block}
\end{column}
\end{columns}
\vspace{0.5em}
\begin{columns}[c]
\begin{column}{0.56\textwidth}
\begin{center}
\alert{Tesis central:} \emph{las tres son empíricamente independientes.}
\end{center}
\vspace{0.2em}
{\footnotesize\color{muted} Un explicador puede ser \textbf{estable y equivocado}, como un reloj parado: siempre marca lo mismo, y siempre mal.}
\end{column}
\begin{column}{0.44\textwidth}\centering
\begin{tikzpicture}[scale=0.9, every node/.style={transform shape}]
  \draw[-{Stealth[length=2mm]},teal,line width=1pt] (0,0) -- (0,1.9) node[above,font=\scriptsize,text=teal]{Estabilidad};
  \draw[-{Stealth[length=2mm]},seafoam,line width=1pt] (0,0) -- (2.3,-0.55) node[right,font=\scriptsize,text=seafoam]{Plausibilidad};
  \draw[-{Stealth[length=2mm]},amber,line width=1pt] (0,0) -- (-2.1,-0.75) node[left,font=\scriptsize,text=amber]{Fidelidad};
  \node[circle,fill=greytx,inner sep=1.3pt] at (0,0) {};
  \node[muted,font=\scriptsize] at (0,-1.5) {tres ejes ortogonales};
\end{tikzpicture}
\end{column}
\end{columns}
\end{frame}

% ##################################################################### 2
\section{Metodología}

\begin{frame}{Un diseño de dos ejes}
\begin{columns}[T]
\begin{column}{0.5\textwidth}
\begin{block}{Eje 1 \textperiodcentered\ Elliptic (real)}
\begin{itemize}\small
  \item Validez \textbf{externa}: transacciones reales de Bitcoin.
  \item Sin \emph{ground-truth} de tipología: solo permite medir \textbf{estabilidad}.
  \item 203.769 nodos, 166 features, partición temporal causal.
  \item Campo receptivo minúsculo ($\sim$2 nodos): solo el ranking de features discrimina.
\end{itemize}
\end{block}
\end{column}
\begin{column}{0.5\textwidth}
\begin{block}{Eje 2 \textperiodcentered\ Sintético (propio)}
\begin{itemize}\small
  \item Validez \textbf{interna}: \emph{ground-truth} por nodo y arista.
  \item Permite medir además \textbf{plausibilidad} y \textbf{fidelidad}.
  \item 4 tipologías canónicas plantadas, con aristas distractoras y firma atenuada.
  \item 3 realizaciones del grafo para verificar.
\end{itemize}
\end{block}
\end{column}
\end{columns}
\vspace{0.3em}
{\footnotesize\color{muted} La validez externa la da Elliptic. La inferencia fuerte sobre plausibilidad y fidelidad la da el sintético, porque Elliptic no dice cuál subgrafo es el patrón.}
\end{frame}

\begin{frame}{Diseño factorial, métricas y protocolo}
\begin{center}
{\color{teal}\Large\textbf{4}} arquitecturas \; $\times$ \; {\color{amber}\Large\textbf{3}} explicadores \; $\times$ \; {\color{seafoam}\Large\textbf{3}} balanceos \; $\times$ \; {\color{coral}\Large\textbf{5}} escenarios
\end{center}
\vspace{0.2em}
\begin{center}\small = 60 configuraciones por eje, cada una replicada con \textbf{3 semillas de modelo}\end{center}
\vspace{0.4em}
\begin{columns}[T]
\begin{column}{0.5\textwidth}
{\footnotesize\color{navy}\textbf{Métricas}}
\begin{itemize}\scriptsize
  \item \textbf{Estabilidad}: Spearman entre rankings de features.
  \item \textbf{Plausibilidad}: coincidencia con el \emph{ground-truth}.
  \item \textbf{Fidelidad}: caída de la predicción al retirar lo importante (Fidelity+).
  \item \textbf{Rendimiento}: PR-AUC, no F1 ni ROC-AUC.
\end{itemize}
\end{column}
\begin{column}{0.5\textwidth}
{\footnotesize\color{navy}\textbf{Protocolo estadístico}}
\begin{itemize}\scriptsize
  \item Pruebas no paramétricas: Kruskal-Wallis y Wilcoxon.
  \item Intervalos de confianza \emph{bootstrap}.
  \item Tamaño de efecto $\eta^2$ de ANOVA.
  \item Corrección de Holm por comparaciones múltiples.
\end{itemize}
\end{column}
\end{columns}
\end{frame}

% ##################################################################### 3
\section{Resultados}

\begin{frame}{Contribución central: dos artefactos de medición}
\begin{columns}[T]
\begin{column}{0.5\textwidth}
\begin{alertblock}{1 \textperiodcentered\ Fallo de memoria silencioso}
\small El cálculo sobre el grafo completo dejaba \textbf{13 configuraciones de GAT} en OOM silencioso, y su media salía solo de las que terminaban. \emph{Favorecía a GAT.}
\end{alertblock}
\end{column}
\begin{column}{0.5\textwidth}
\begin{alertblock}{2 \textperiodcentered\ Truncamiento del Spearman}
\small La métrica dimensionaba los rangos por \texttt{top\_k} y descartaba las features fuera del umbral, mutilando el ranking. \emph{Favorecía a GraphSAGE.}
\end{alertblock}
\end{column}
\end{columns}
\vspace{0.5em}
\begin{center}\small
\alert{Cada artefacto, por separado, bastaba para una conclusión comparativa falsa.}\\ La estabilidad medida depende del protocolo, no solo del método.
\end{center}
\vspace{0.2em}
{\scriptsize\color{muted}\hfill Además, dos \emph{bugs} reportables en el PGExplainer de PyG 2.7. En línea con Kosan et al. (2023).\hspace*{0.3cm}}
\end{frame}

% -------- Resultado 1: figura conceptual de particion (reemplaza barras)
\begin{frame}{Resultado 1: la estabilidad separa dos grupos, no cuatro puestos}
\begin{columns}[T]
\begin{column}{0.42\textwidth}
\footnotesize
La corrección de la métrica, con la matriz reentrenada a \textbf{3 semillas (180 modelos)}, no revela un ranking de cuatro, sino una \alert{partición en dos grupos}.
\vspace{0.4em}

La línea \textbf{no} se traza por la distancia entre barras, sino por dónde la diferencia es \textbf{consistente}:
\begin{itemize}\scriptsize
  \item \textbf{Entre} grupos: significativa.
  \item \textbf{Dentro} de cada grupo: no (Wilcoxon pareado $p=0{,}17$ y $p=0{,}10$).
\end{itemize}
\vspace{0.2em}
{\scriptsize\color{muted} GAT y GCN se permutan entre semillas: no hay un ganador único.}
\end{column}
\begin{column}{0.58\textwidth}\centering
\begin{tikzpicture}[scale=0.55]
  % marco de ejes: eje x = estabilidad de 0.64 a 0.82 -> 0 a 12
  % transform: x(v) = (v-0.64)*66.667
  \def\Y{2.9}
  % bandas de grupo
  \begin{scope}[on background layer]
    \fill[muted!14,rounded corners=2pt] (-0.3,-0.15) rectangle (7.05,\Y+0.35);
    \fill[seafoam!16,rounded corners=2pt] (7.25,-0.15) rectangle (12.4,\Y+0.35);
  \end{scope}
  % eje x
  \draw[muted,line width=0.7pt] (-0.3,-0.25) -- (12.4,-0.25);
  \foreach \v/\lab in {0.65/0{,}65, 0.70/0{,}70, 0.75/0{,}75, 0.80/0{,}80}{
    \pgfmathsetmacro\xx{(\v-0.64)*66.667}
    \draw[muted,line width=0.6pt] (\xx,-0.25) -- (\xx,-0.4);
    \node[muted,font=\tiny] at (\xx,-0.62) {\lab};
  }
  \node[muted,font=\scriptsize] at (6.05,-1.0) {estabilidad (Spearman de features)};
  % divisor / brecha
  \draw[coral,dashed,line width=1pt] (7.15,-0.15) -- (7.15,\Y+0.35);
  \node[coral,font=\scriptsize,align=center] at (7.15,\Y+0.75) {brecha\\significativa};
  % etiquetas de grupo
  \node[muted,font=\scriptsize] at (2.2,\Y+0.55) {grupo bajo};
  \node[seafoam,font=\scriptsize] at (9.8,\Y+0.55) {grupo alto};
  % puntos con barras de error horizontales
  % GAT 0.782 +-0.013 -> x=9.47 err 0.867 ; y alto
  % GCN 0.758 +-0.025 -> x=7.87 err 1.667
  % SAGE 0.735 +-0.022 -> x=6.33 err 1.467
  % TAGCN 0.672 +-0.077 -> x=2.13 err 5.133
  \ptmark{9.47}{0.87}{2.35}{seafoam}{GAT}{0{,}782}
  \ptmark{7.87}{1.67}{1.55}{seafoam}{GCN}{0{,}758}
  \ptmark{6.33}{1.47}{0.78}{muted}{GraphSAGE}{0{,}735}
  \ptmark{2.13}{5.13}{0.10}{muted}{TAGCN}{0{,}672}
\end{tikzpicture}
\figcap{Estabilidad por arquitectura con IC 95\%. Los grupos se solapan por dentro y se separan en la brecha. TAGCN tiene la mayor dispersión. \emph{Fuente: elaboración propia.}}
\end{column}
\end{columns}
\end{frame}

\begin{frame}{Resultado 1b: la partición concuerda en ambos regímenes}
\begin{columns}[T]
\begin{column}{0.5\textwidth}
\begin{itemize}\small
  \item La \emph{inversión por densidad} que reportaba una versión anterior era un \textbf{artefacto} de la métrica.
  \item Corregida, Elliptic (disperso) \textbf{concuerda} con el sintético (denso), sobre datos y \emph{pipeline} independientes.
  \item El sintético separa \textbf{los mismos dos grupos}.
\end{itemize}
\vspace{0.4em}
\begin{center}
\small Correlación de rangos entre regímenes\\[3pt]
{\color{coral}\Large\textbf{$-0{,}20$}} \; $\rightarrow$ \; {\color{amber}\Large\textbf{$+0{,}80$}}\\[2pt]
{\footnotesize métrica con \emph{bug} \; a \; métrica corregida}
\end{center}
\end{column}
\begin{column}{0.5\textwidth}\centering
\figcard{chapter_5/images_ch5/contraste_regimen.png}
\figcap{Orden de estabilidad: Elliptic (disperso) frente al sintético (denso). \emph{Fuente: elaboración propia.}}
\end{column}
\end{columns}
\end{frame}

\begin{frame}{Resultado 2: el más plausible no es el más fiel}
\begin{columns}[T]
\begin{column}{0.52\textwidth}
\begin{itemize}\small
  \item \textbf{\color{teal}PGExplainer} domina la plausibilidad de aristas: 0,80 frente a 0,50, con 0,40 de azar (Wilcoxon $p\approx 2{,}6\times10^{-35}$).
  \item Pero \textbf{\color{coral}colapsa en fidelidad}: 0,11 frente a 0,56 de GNNExplainer.
  \item \emph{El explicador más plausible no es el más fiel.} Confirma la independencia de las dimensiones.
\end{itemize}
\vspace{0.3em}
{\scriptsize\color{muted} La figura compara los dos explicadores con máscara de aristas. GNNShap no la produce por diseño: su plausibilidad es de features, con otra línea base, y va en el respaldo.}
\end{column}
\begin{column}{0.48\textwidth}\centering
\figcard{disociacion_eb.png}
\figcap{Plausibilidad de aristas frente a Fidelity+ (eje sintético, IC 95\%, azar de aristas $=0{,}40$). \emph{Fuente: elaboración propia.}}
\end{column}
\end{columns}
\end{frame}

\begin{frame}{Resultado 3: el puente que no existe}
\begin{columns}[T]
\begin{column}{0.52\textwidth}
\begin{itemize}\small
  \item Nuestra hipótesis central (H3): una explicación más estable señalaría mejor el patrón real.
  \item El dato dice que \textbf{no}. Correlación entre estabilidad y plausibilidad $r=-0{,}01$, con IC 95\,\% de $-0{,}038$ a $+0{,}011$.
  \item El intervalo \textbf{incluye el cero}: no hay relación.
\end{itemize}
\vspace{0.4em}
\small\emph{Es un resultado nulo, y lo reportamos porque contradice lo que esperábamos.}
\end{column}
\begin{column}{0.48\textwidth}\centering
\figcard{chapter_5/images_ch5/puente_nulo.png}
\figcap{Relación entre estabilidad y plausibilidad (GNNExplainer): $r=-0{,}01$, IC incluye 0. \emph{Fuente: elaboración propia.}}
\end{column}
\end{columns}
\end{frame}

% -------- Resultado 4: figura conceptual esperado vs observado (reemplaza barras)
\begin{frame}{Resultado 4: el desbalance no gobierna la estabilidad}
\begin{columns}[T]
\begin{column}{0.42\textwidth}
\footnotesize
Esperábamos (H1) que la estabilidad \textbf{cayera} al agravarse el desbalance. La curva observada es \alert{plana y sin dirección}.
\begin{itemize}\scriptsize
  \item Rango de 0,696 (1:1) a 0,765 (1:50), \textbf{siete centésimas}, menor que las once entre arquitecturas.
  \item \textbf{Kruskal-Wallis $p=0{,}18$}: no se rechaza la igualdad. Efecto $\eta^2=0{,}05$ (pequeño).
  \item El mínimo está en \textbf{1:1}, el más equilibrado. La variación no tiene dirección.
\end{itemize}
\begin{block}{}
\scriptsize Implicación: el balanceo puede elegirse por \textbf{rendimiento}, sin temer que degrade la interpretabilidad.
\end{block}
\end{column}
\begin{column}{0.58\textwidth}\centering
\begin{tikzpicture}[scale=0.62]
  % eje y estabilidad 0.5..0.85 -> 0..6 ; y(v)=(v-0.5)*17.143 (valores numericos)
  \draw[muted,line width=0.7pt] (-0.2,0) -- (-0.2,6.1);
  \draw[muted,line width=0.7pt] (-0.2,0) -- (11.2,0);
  \foreach \yy/\v in {1.714/0.6, 3.429/0.7, 5.143/0.8}{
    \draw[muted,line width=0.5pt] (-0.35,\yy) -- (-0.2,\yy);
    \node[muted,font=\tiny,anchor=east] at (-0.4,\yy) {\v};
  }
  \node[muted,font=\scriptsize,rotate=90] at (-1.15,3.0) {estabilidad};
  \foreach \x/\lab in {0/1:1, 2.6/1:10, 5.2/1:50, 7.8/1:100, 10.4/nativo}{
    \node[muted,font=\tiny,anchor=north] at (\x,-0.1) {\lab};
  }
  \node[muted,font=\scriptsize] at (5.2,-0.85) {escenario de desbalance};
  % linea esperada H1 (cae): valores numericos ya evaluados
  \draw[coral,dashed,line width=1.2pt]
     (0,5.143) -- (2.6,3.771) -- (5.2,2.400) -- (7.8,1.200) -- (10.4,1.714);
  \node[coral,font=\scriptsize,anchor=west] at (7.9,1.55) {lo que H1 predecía};
  % linea observada (plana): 0.696,0.736,0.765,0.745,0.760
  \draw[amber,line width=1.4pt]
     (0,3.360) -- (2.6,4.046) -- (5.2,4.543) -- (7.8,4.200) -- (10.4,4.457);
  \foreach \x/\y in {0/3.360, 2.6/4.046, 5.2/4.543, 7.8/4.200, 10.4/4.457}{
    \node[circle,fill=amber,inner sep=2pt] at (\x,\y) {};
  }
  \node[amber,font=\scriptsize,anchor=south] at (5.2,4.75) {observado ($p=0{,}18$)};
  % marca del minimo en 1:1
  \draw[seafoam,line width=0.8pt] (0,3.360) circle (0.22);
  \node[seafoam,font=\tiny,anchor=south west] at (0.2,2.85) {mínimo en 1:1};
\end{tikzpicture}
\figcap{Estabilidad por escenario: la predicción de H1 frente a la curva observada, plana. \emph{Fuente: elaboración propia.}}
\end{column}
\end{columns}
\end{frame}

\begin{frame}{Rendimiento: colapso validación a test, por eso PR-AUC}
\begin{columns}[c]
\begin{column}{0.56\textwidth}\centering
\figcard[\linewidth]{colapso_val_test.png}
\figcap{PR-AUC y precisión@50 colapsan de validación a test, mientras el ROC-AUC cae mucho menos y aparenta un modelo aceptable. \emph{Fuente: elaboración propia.}}
\end{column}
\begin{column}{0.44\textwidth}
\small
\begin{itemize}
  \item El modelo aprende en validación y \textbf{colapsa en test}: es el \emph{shift} temporal del dataset, una propiedad del dato, no del método.
  \item El \textbf{ROC-AUC engaña} bajo desbalance (0,88). Tomamos \textbf{PR-AUC} y \textbf{precisión@k} como primarias.
  \item Por eso la estabilidad se estudia sobre los \textbf{verdaderos positivos de validación}, encuadre declarado en el manuscrito.
\end{itemize}
\end{column}
\end{columns}
\end{frame}

% ##################################################################### 4
\section{Conclusiones}

\begin{frame}{Las tres hipótesis, y qué pasó con cada una}
\begin{center}
\small
\tabcap{Hipótesis de trabajo y su veredicto empírico.}
\renewcommand{\arraystretch}{1.5}
\begin{tabular}{p{0.42\textwidth}p{0.15\textwidth}p{0.33\textwidth}}
\toprule
\textbf{Lo que predijimos} & \textbf{Veredicto} & \textbf{La evidencia} \\
\midrule
\textbf{H1.} La estabilidad se degradaría con el desbalance. & \textcolor{coral}{\textbf{Refutada}} & Variación sin dirección: $p=0{,}18$, mínimo en 1:1. \\
\textbf{H2.} TAGCN sería la más estable. & \textcolor{coral}{\textbf{Refutada}} & TAGCN en el grupo bajo en las 3 semillas. \\
\textbf{H3.} Un explicador más estable señalaría mejor el patrón. & \textcolor{coral}{\textbf{Refutada}} & Puente nulo: $r=-0{,}01$, IC incluye el cero. \\
\bottomrule
\end{tabular}
\end{center}
\vspace{0.5em}
{\footnotesize Las tres se cayeron y las tres se reportan. Escribir las predicciones antes de ver los datos es lo que convierte su refutación en un resultado.}
\end{frame}

\begin{frame}{Conclusiones y matriz de recomendación}
\begin{columns}[T]
\begin{column}{0.54\textwidth}
{\footnotesize\color{navy}\textbf{Conclusiones}}
\begin{enumerate}\scriptsize
  \item \textbf{Tres dimensiones independientes}: no existe una única buena explicación.
  \item \textbf{Dos grupos de estabilidad}, no cuatro puestos, y la partición concuerda entre datos reales y sintéticos.
  \item \textbf{El explicador es la palanca dominante} ($\eta^2$ de 0,36 a 0,64). Balanceo despreciable ($\eta^2\le 0{,}01$), escenario pequeño y no significativo.
  \item \textbf{El más plausible no es el más fiel}.
  \item \textbf{Aporte metodológico}: dos artefactos corregidos y reproducibilidad de pesos frente a la de conclusiones.
\end{enumerate}
\end{column}
\begin{column}{0.46\textwidth}
\tabcap{Recomendación por objetivo de auditoría.}
{\footnotesize
\renewcommand{\arraystretch}{1.3}
\begin{tabular}{ll}
\toprule
\textbf{Objetivo} & \textbf{Configuración} \\
\midrule
Auditabilidad & \textcolor{teal}{grupo alto: GAT o GCN} \\
Recuperar patrón & \textcolor{seafoam}{PGExplainer} \\
Fidelidad al modelo & \textcolor{amber}{GNNExplainer} \\
Estabilidad interna & \textcolor{coral}{GNNShap} \\
\bottomrule
\end{tabular}}
\vspace{0.3em}
{\scriptsize\color{muted} La recomendación de arquitectura es de \textbf{grupo}: dentro del alto la evidencia no distingue, así que se elige por coste o disponibilidad.}
\end{column}
\end{columns}
\end{frame}

\begin{frame}[standout]
Gracias \textperiodcentered\ Preguntas
\vspace{0.5em}

{\small\color{greytx} Estabilidad $\neq$ plausibilidad $\neq$ fidelidad: el mejor explicador depende del objetivo.\\
Trabajo futuro: Elliptic2, AMLSim, GNNs temporales, GraphSMOTE.}
\vspace{0.5em}

{\footnotesize\color{muted} Alejandro Gómez Huertas \textperiodcentered\ Juan Diego Garzón Ovalle}
\end{frame}

% ##################################################################### RESPALDO
\appendix

\begin{frame}[allowframebreaks]{Referencias seleccionadas}
\scriptsize
\begin{itemize}
  \item Weber M y colaboradores (2019). \emph{Anti-Money Laundering in Bitcoin: Experimenting with Graph Convolutional Networks for Financial Forensics.} arXiv:1908.02591.
  \item Kipf T y Welling M (2017). \emph{Semi-Supervised Classification with Graph Convolutional Networks.} ICLR.
  \item Hamilton W, Ying R y Leskovec J (2017). \emph{Inductive Representation Learning on Large Graphs.} NeurIPS.
  \item Veličković P y colaboradores (2018). \emph{Graph Attention Networks.} ICLR.
  \item Du J y colaboradores (2017). \emph{Topology Adaptive Graph Convolutional Networks.} arXiv:1710.10370.
  \item Ying R y colaboradores (2019). \emph{GNNExplainer: Generating Explanations for Graph Neural Networks.} NeurIPS.
  \item Luo D y colaboradores (2020). \emph{Parameterized Explainer for Graph Neural Network.} NeurIPS.
  \item Akkas S y Azad A (2024). \emph{GNNShap: Scalable and Accurate GNN Explanation using Shapley Values.} WWW.
  \item Kosan M y colaboradores (2023). \emph{GNNX-BENCH: Unravelling the Utility of Perturbation-based GNN Explainers through In-depth Benchmarking.} arXiv:2310.01794.
  \item Agarwal C, Zitnik M y Lakkaraju H (2022). \emph{Probing GNN Explainers: A Rigorous Theoretical and Empirical Analysis of GNN Explanation Methods.} AISTATS.
  \item He X y colaboradores (2026). \emph{An Explainable Graph Neural Network Framework for Illicit Financial Transaction Detection.} Applied Intelligence 56.
  \item Zhao T, Zhang X y Wang S (2021). \emph{GraphSMOTE: Imbalanced Node Classification on Graphs with Graph Neural Networks.} WSDM.
\end{itemize}
\end{frame}

\begin{frame}{Respaldo \textperiodcentered\ estabilidad por semilla de modelo}
\begin{center}
\small
\tabcap{Estabilidad por semilla de modelo (Elliptic, GNNExplainer).}
\renewcommand{\arraystretch}{1.25}
\begin{tabular}{lccccc}
\toprule
\textbf{Arquitectura} & \textbf{s42} & \textbf{s43} & \textbf{s44} & \textbf{Media} & \textbf{IC95\,\%} \\
\midrule
GAT       & 0,766 & 0,789 & 0,790 & \textbf{0,782} & 0,758 a 0,805 \\
GCN       & 0,783 & 0,733 & 0,757 & \textbf{0,758} & 0,724 a 0,787 \\
GraphSAGE & 0,756 & 0,713 & 0,738 & 0,735 & 0,710 a 0,756 \\
TAGCN     & 0,586 & 0,693 & 0,736 & 0,672 & 0,632 a 0,717 \\
\bottomrule
\end{tabular}
\end{center}
\vspace{0.5em}
{\footnotesize \textbf{GAT y GCN se permutan} entre semillas: por eso no nombramos un ganador único. \textbf{TAGCN} tiene la mayor dispersión (dt 0,077, el triple): su lugar en el grupo bajo es firme, su valor central admite margen.}
\end{frame}

\begin{frame}{Respaldo \textperiodcentered\ disociación: valores exactos}
\begin{center}
\small
\tabcap{Plausibilidad y Fidelity+ por explicador (eje sintético, g0, 3 semillas). Cada explicador se compara contra la línea base de azar de su propia métrica.}
\renewcommand{\arraystretch}{1.3}
\begin{tabular}{lccc}
\toprule
\textbf{Explicador} & \textbf{Plaus. aristas} & \textbf{Plaus. features} & \textbf{Fidelity+} \\
\midrule
GNNExplainer & 0,50 & \textbf{0,45} & \textbf{0,56} \\
PGExplainer & \textbf{0,80} & n/d & 0,11 \\
GNNShap & n/d & 0,15 & 0,46 \\
\midrule
\emph{Azar (línea base)} & \emph{0,40} & \emph{0,08} & no aplica \\
\bottomrule
\end{tabular}
\end{center}
\vspace{0.3em}
{\scriptsize n/d = métricas que el explicador no produce por diseño. Contra su propia línea base, PGExplainer duplica el azar en aristas y GNNExplainer lo sextuplica en features. \textbf{GNNShap} es el más estable internamente (0,98).}
\end{frame}

\begin{frame}{Respaldo \textperiodcentered\ ¿por qué GNNExplainer para el ranking?}
\begin{center}
\small
\tabcap{Estabilidad (Spearman) por explicador y arquitectura en Elliptic.}
\renewcommand{\arraystretch}{1.25}
\begin{tabular}{lcccc}
\toprule
\textbf{Explicador (Elliptic)} & \textbf{GAT} & \textbf{GCN} & \textbf{GraphSAGE} & \textbf{TAGCN} \\
\midrule
\textbf{\color{teal}GNNExplainer} & 0,78 & 0,76 & 0,74 & 0,67 \\
PGExplainer & 0,00 & 0,00 & 0,00 & 0,00 \\
GNNShap & 0,97$^{*}$ & 0,97 & 0,95 & 0,97$^{*}$ \\
\bottomrule
\end{tabular}
\end{center}
\vspace{0.5em}
{\footnotesize \textbf{\color{teal}GNNExplainer} es el único que discrimina entre arquitecturas: PGExplainer degenera (Spearman nulo) y GNNShap se satura cerca de 0,96. Es además el explicador común a los dos ejes. GNNShap a 3 semillas cubre GCN y GraphSAGE. Las marcadas con $^{*}$, GAT y TAGCN, solo en la semilla 42.}
\end{frame}

\begin{frame}{Respaldo \textperiodcentered\ curvas ROC y precisión-recall}
\begin{center}
\figcard[0.82\linewidth]{curvas_pr_roc.png}
\end{center}
\figcap{Curvas ROC (izquierda) y precisión-recall (derecha), validación frente a test, mismos modelos. El ROC cae de 0,88 a 0,65, pero la precisión-recall se desploma casi a la tasa base de ilícitas. \emph{Fuente: elaboración propia.}}
\end{frame}

\end{document}
